\documentclass[10pt,a4paper]{amsart}
\usepackage[margin=19mm]{geometry}
\usepackage{amsmath,amssymb,mathtools}
\usepackage[scr=rsfs,cal=euler]{mathalfa}
\usepackage{xfrac}
\usepackage[unicode,hidelinks,bookmarks=false]{hyperref}
\newcommand{\ie}{i.e.\ }
\newcommand{\cf}{cf.\ }
\newcommand{\resp}{respectively\ }
\newcommand{\Subsection}{\subsection}
\providecommand{\nobreakdash}{\nobreak\hbox{-}}
\setlength{\emergencystretch}{2em}
\multlinegap0pt
\usepackage{bm}
\usepackage{bbm}
\usepackage{dsfont}
\usepackage{xspace}
\usepackage{cancel}
\usepackage{mathtools}
\usepackage{amsthm}
\makeatletter
\@ifundefined{lemma}{\newtheorem{lemma}{Lemma}}{}
\makeatother
\newcommand{\Fone}{\mathcal{F}_Z}
\newcommand{\Ftwo}{\mathcal{F}_{Z,NL}}
\newcommand{\NN}{\mathcal{N}}
\newcommand{\R}{\mathbb R}
\newcommand{\brc}[1]{\left\{#1\right\}}
\newcommand{\y}{\mathbf{y}}
\newcommand{\zero}{\mathbf{0}}
\title[Nonlinear Network Identifiability with Full Excitations]{Nonlinear Network Identifiability with Full Excitations}
\author{Renato Vizuete, Julien M. Hendrickx}
\date{Source: arXiv:2405.07636v3 (CC BY 4.0)}
\begin{document}
\maketitle

\noindent\emph{Source and licence.} Nonlinear Network Identifiability with Full Excitations. Version arXiv:2405.07636v3 (CC BY 4.0), distributed under the Creative Commons Attribution 4.0 International licence, as recorded in the arXiv licence metadata for this exact version and shown on the abstract page https://arxiv.org/abs/2405.07636v3. Full rights notice: licenses/arxiv-2405.07636-CC-BY-4.0.txt.\par\smallskip

\noindent\textit{Editorial scope.} Counted unit: the Lemma whose source label is \texttt{lemma:recursively} in the article (source0.tex lines 751--754, source label \texttt{lemma:recursively}), delivered together with the article's own complete original proof of it (source0.tex lines 755--790, one \texttt{proof} environment of 36 lines). No mathematics is written, repaired or re-derived: statement and proof are the article's own text line for line. Required context retained uncounted: lemma:sinks (lines 445--447); lemma:unique\_path (lines 630--632); lemma:removal\_node (lines 656--659); lemma:multivariate\_sum (lines 735--747). Every definition, display and label of those lines is retained unchanged. Omitted from this package and not redistributed: 1--444, 448--629, 633--655, 660--734, 748--750, 791--1291. Nothing inside a retained excerpt is omitted or altered: every retained body is delivered exactly as the article's lines are written, so no omission receipt of any kind accompanies this package. Citations: the retained excerpts carry no citation command at all, so no bibliography entry of the article is transcribed and no reference is redistributed. Reference adaptation: none was needed and none was made; every reference inside the delivered unit resolves inside it, so no unresolved reference or citation is produced. Printed slips kept literal: no printed word, symbol, hypothesis or number of the counted statement or of its proof was altered. Layout and class: the article's own class is \texttt{IEEEtran} with packages including graphics, epsfig, amsmath, amssymb, dsfont, enumerate, subfigure, multirow, bm, tikz, cite, amsthm, hyperref, algorithm2e; the wrapper is the lane's pinned \texttt{amsart} editorial wrapper at 10pt on A4, which restates the theorem-like environments the retained text needs and adds the article's own macro definitions that the retained text needs (\texttt{\textbackslash Fone}, \texttt{\textbackslash Ftwo}, \texttt{\textbackslash NN}, \texttt{\textbackslash R}, \texttt{\textbackslash brc}, \texttt{\textbackslash y}, \texttt{\textbackslash zero}), with extra wrapper packages (bm, bbm, dsfont, xspace, cancel, mathtools). The delivered numbering is the unit's own. Assets: no retained span contains a figure environment, a graphics-inclusion command, a PDF-page inclusion, a table or a TikZ picture; no image file is copied into this package, the package declares no asset dependency and no image file is redistributed with it. Licence: version arXiv:2405.07636v3 of this article is distributed under the Creative Commons Attribution 4.0 International licence, as recorded in the arXiv licence metadata for this exact version and shown on the abstract page https://arxiv.org/abs/2405.07636v3. A full rights notice is delivered at licenses/arxiv-2405.07636-CC-BY-4.0.txt. Characters: every retained excerpt is pure ASCII, so the delivered engine, XeLaTeX with its default Latin Modern text font, reports no missing character.
\begin{lemma}\label{lemma:sinks}
Let $G$ be a DAG such that the edge $(i,j)$ is the only path from the node $j$ to the node $i$. Then if $i$ is measured, $f_{i,j}$ is identifiable in the class $\Fone$.
\end{lemma}

\begin{lemma}\label{lemma:unique_path}
    In a DAG, any node $i$ which is not a source has at least one in-neighbor $j$ with only one path to $i$.   
\end{lemma}

\begin{lemma}[Removal of a node]\label{lemma:removal_node}
    Given a DAG $G$ and a node $j$ with only one path to its out-neighbor $i$. Let us assume that $j$ has been measured and the edge $f_{i,j}$ is identifiable. When $i\in \NN^m$, the edges $f_{i,\ell}$ and functions $F_\ell$ for $\ell\in\NN_i$ are identifiable in $G$ if they are identifiable in the induced subgraph   
    $G_{V\setminus\{ j\}}$.
\end{lemma}

\begin{lemma}[Sum of multivariate functions]\label{lemma:multivariate_sum}
Given three non identically zero twice continuously differentiable functions $f:\R^p\to \R$ and $g,\tilde g:\R^q\to\R$. Let us assume that at least one of the partial derivatives of $f$ is not constant and they satisfy:
$$
f(\zero)= g(\zero)= \tilde g(\zero)=0;
$$
\vspace{-7mm}
\begin{multline}\label{eq:statement_lemma_function_h}
f(x_1\!+\!g(\y_1,\ldots,\y_r),\ldots,x_p\!+\!g(\tau^{p-1}\y_1,\ldots,\tau^{p-1}\y_r))=\\     
f(x_1\!+\!\tilde g(\y_1,\ldots,\y_r),\ldots,x_p\!+\!\tilde g(\tau^{p-1}\y_1,\ldots,\tau^{p-1}\y_r))+h,
\end{multline}
for all $x_i$ with $i\in \{1,\ldots,p\}$, and for all $\y_j$ with $j\in \{1,\ldots,r\}$, where $h$ is an arbitrary function that does not depend on $x_1,\ldots,x_p$.
Then $g=\tilde g$.
\end{lemma}

\begin{lemma}\label{lemma:recursively}
 In a DAG and in the class $\Ftwo$, if the function $F_i$ is identifiable, then all the functions $f_{i,j}$ and 
$F_j$ are identifiable for all $j\in \NN_i$.
\end{lemma}

\begin{proof}
    If $F_i$ corresponds to a source, the result is trivial. Let us assume that $F_i$ corresponds to a node that is not a source. 
    The output of any node $i$ is of the form:
    \begin{align}
        y_i^k&=u_i^{k-1}+F_i\nonumber\\
        &=u_i^{k-1}\!+\!\sum_{j\in\NN_i}\!f_{i,j}(u_{j}^{k-2}\!+\!F_{j}^{(1)},\ldots,u_{j}^{k-m_{j}-1}\!+\!F_{j}^{(m_{j})}).\nonumber
    \end{align}
    Let us assume that
    there exists a set $\{\tilde f\}\neq \{ f\}$ such that $F_i=\tilde F_i$, which implies:
\begin{multline}\label{eq:equal_f_tilde_DAG}
        \sum_{j\in\NN_i}f_{i,j}(u_{j}^{k-2}+F_{j}^{(1)},\ldots,u_{j}^{k-m_{j}-1}+F_{j}^{(m_{j})})=\\
        \sum_{j\in\NN_i}\tilde f_{i,j}(u_{j}^{k-2}+\tilde F_{j}^{(1)},\ldots,u_{j}^{k-m_{j}-1}+\tilde F_{j}^{(m_{j})}).
    \end{multline}
By Lemma~\ref{lemma:unique_path}, there is at least one in-neighbor $\ell$ with only one path to $i$. 
Then, we can use Lemma~\ref{lemma:sinks} to guarantee $f_{i,\ell}=\tilde f_{i,\ell}$, and \eqref{eq:equal_f_tilde_DAG} can be expressed as:
\begin{multline}\label{eq:fij_lemma5}
    f_{i,\ell}(u_{\ell}^{k-2}+F_{\ell}^{(1)},\ldots,u_{\ell}^{k-m_{\ell}-1}+F_{\ell}^{(m_{\ell})})=\\
         f_{i,\ell}(u_{\ell}^{k-2}+\tilde F_{\ell}^{(1)},\ldots,u_{\ell}^{k-m_{\ell}-1}+\tilde F_{\ell}^{(m_{\ell})})+\phi,
\end{multline}
where $\phi$ does not depend on any of the inputs $u_\ell$ since the only path from $\ell$ to $i$ is the edge $(i,\ell)$. 
By using Lemma~\ref{lemma:multivariate_sum} in \eqref{eq:fij_lemma5} we obtain $F_{\ell}^{(1)}=\tilde F_{\ell}^{(1)}$, which corresponds to the mathematical constraint associated with the measurement of the node $\ell$.
Then, the equality \eqref{eq:equal_f_tilde_DAG} becomes:
\begin{multline}\label{eq:node_i2_Volterra}
\sum_{j\in\NN_i\setminus\brc{\ell}}f_{i,j}(u_{j}^{k-2}+F_{j}^{(1)},\ldots,u_{j}^{k-m_{j}-1}+F_{j}^{(m_{j})})=\\
\sum_{j\in\NN_i\setminus\brc{\ell}}\tilde f_{i,j}(u_{j}^{k-2}+\tilde F_{j}^{(1)},\ldots,u_{j}^{k-m_{j}-1}+\tilde F_{j}^{(m_{j})}),
\end{multline}
and according to Lemma~\ref{lemma:removal_node}, the edges $f_{i,j}$ and functions $F_j$ for $j\in\NN_i\setminus\{\ell\}$ are identifiable if they are identifiable in the induced subgraph $G_{V\setminus\{ \ell\}}$.
By Lemma~\ref{lemma:unique_path}, in the subgraph $G_{V\setminus\{ \ell\}}$ there must be now another node with only one path to the node $i$ and we can follow a similar approach. Then, we can show that 
$$
f_{i,j}=\tilde f_{i,j} \quad \text{ for all } j\in\NN_i,
$$
and similarly for the functions
$$
F_j^{(1)}=\tilde F_j^{(1)} \quad \text{ for all } j\in\NN_i.
$$
\end{proof}

\end{document}
